\chapter{Heat Rash (Miliaria)}

\section*{Definition}
A skin condition caused by blockage of eccrine sweat ducts during hot, humid weather, leading to sweat retention and superficial inflammation presenting as small red bumps or blisters.

\section*{What Happens in Your Body}
Sweat duct obstruction from keratin plugs, bacteria, or inflammation leads to sweat retention and ductal rupture. Miliaria crystallina (superficial, clear vesicles) involves stratum corneum duct obstruction. Miliaria rubra (prickly heat, red papules) involves deeper epidermal duct rupture with periductal inflammation. Miliaria profunda (deep, flesh-colored papules) involves dermal duct rupture.

\section*{Causes}
\begin{itemize}
  \item Hot, humid weather (most common)
  \item Excessive sweating (exercise, fever)
  \item Tight or non-breathable clothing
  \item Prolonged bed rest (friction and heat trapping)
  \item Occlusive dressings or ointments
  \item Neonatal (immature sweat ducts in infants)
  \item Fever with profuse sweating
  \item Intense physical activity in heat
  \item Overdressing in warm environments
  \item Tropical climate exposure
\end{itemize}

\section*{When to See a Doctor}
See your doctor if:
\begin{itemize}
  \item Symptoms are severe or getting worse over time
  \item Heat rash with fever, signs of heat exhaustion (nausea, dizziness, confusion), or rash with pus drainage suggesting secondary infection
  \item You have other concerning symptoms such as fever, weight loss, or bleeding
\end{itemize}

\section*{Related Symptoms}
\begin{itemize}
  \item Itching (pruritus)
  \item Red bumps on skin
  \item Prickling sensation
  \item Heat exhaustion
  \item Dehydration
\end{itemize}
\textbf{Important:} If this symptom persists, your doctor will check for serious underlying causes.

\section*{Self-Care / Home Management}
\begin{itemize}
  \item Move to a cool, shaded, or air-conditioned environment and allow the skin to dry without occlusion.
  \item Wear loose, lightweight, breathable fabrics (cotton or moisture-wicking materials) to reduce sweating and friction.
  \item Apply cool compresses or take cool showers to lower skin temperature and soothe irritation.
  \item Use calamine lotion or over-the-counter 1\% hydrocortisone cream sparingly for itching; avoid heavy ointments that block sweat ducts.
\end{itemize}

\section*{How Your Doctor Will Evaluate You}
\begin{itemize}
  \item Visual inspection of the rash — distribution, morphology (clear vesicles vs red papules vs deep flesh-coloured bumps), and location in sweat-prone areas — is usually sufficient for diagnosis.
  \item History of recent heat exposure, sweating, occlusive clothing, or fever helps differentiate miliaria from other skin conditions.
  \item No laboratory or imaging tests are routinely needed; skin biopsy is reserved for atypical or persistent cases where diagnosis is unclear.
  \item Referral to a dermatologist is arranged if the rash does not resolve with cooling measures or if secondary infection is suspected.
\end{itemize}

\section*{Treatment Options}
\begin{itemize}
  \item The primary treatment is removing the aggravating factor — move to a cool environment, remove occlusive clothing, and allow the skin to dry and breathe.
  \item Topical calamine lotion or 1\% hydrocortisone cream can reduce inflammation and itching; oral antihistamines may help if itching is significant.
  \item Miliaria profunda may require topical keratolytic agents (such as salicylic acid or lactic acid) to unblock sweat ducts under dermatological guidance.
  \item Secondary bacterial infection (impetiginisation) should be treated with topical or oral antibiotics as prescribed; avoid using heavy emollients or ointments.
  \item Medications, lifestyle modifications, and physical therapies may be prescribed as appropriate.
  \item Follow-up ensures treatment effectiveness and allows timely adjustments to the care plan.
\end{itemize}

\section*{References}

\begin{thebibliography}{99}
\bibitem{harrison-derm} Longo DL, et al. \emph{Harrison's Principles of Internal Medicine}. 21st ed. McGraw-Hill; 2022. Chapter 57: Miliaria.
\bibitem{uptodate-heatrash} Miller JL, et al. Miliaria. In: UpToDate, Post TW (Ed), UpToDate, Waltham, MA; 2025.
\bibitem{guerrero} Guerrero V, et al. Heat-related skin disorders. \emph{Dermatol Clin}. 2023;41(3):421-432.
\bibitem{booth} Booth B, et al. Miliaria: a clinical update. \emph{Aust Fam Physician}. 2022;51(5):310-314.
\bibitem{fitzpatrick} Kang S, et al. \emph{Fitzpatrick's Dermatology}. 10th ed. McGraw-Hill; 2023. Chapter 89: Disorders of Eccrine Sweat Glands.
\bibitem{bolognia} Bolognia JL, et al. \emph{Dermatology}. 4th ed. Elsevier; 2022. Chapter 39: Miliaria.
\end{thebibliography}
